
\documentclass{article}
\usepackage{colortbl}
\usepackage{makecell}
\usepackage{multirow}
\usepackage{supertabular}

\begin{document}

\newcounter{utterance}

\centering \large Interaction Transcript for game `ifeval', experiment `ifeval', episode 339 with Qwen3.5{-}4B.
\vspace{24pt}

{ \footnotesize  \setcounter{utterance}{1}
\setlength{\tabcolsep}{0pt}
\begin{supertabular}{c@{$\;$}|p{.15\linewidth}@{}p{.15\linewidth}p{.15\linewidth}p{.15\linewidth}p{.15\linewidth}p{.15\linewidth}}
    \# & \multicolumn{2}{c}{Player} && \multicolumn{2}{c}{Game Master} \\
    \hline

    \theutterance \stepcounter{utterance}  
    & & & \multicolumn{4}{p{0.6\linewidth}}{
        \cellcolor[rgb]{0.9,0.9,0.9}{
            \makecell[{{p{\linewidth}}}]{
                \texttt{\tiny{[P1$\langle$GM]}}
                \texttt{What is inside Shinto shrines? Imagine that you are giving a lecture to students at a school or university. Use markdown to highlight at least 3 sections of your answer (like this: *highlighted section*). Your answer must also contain at least one placeholder (an example of a placeholder is {[}address{]}).} \\
            }
        }
    }
     \\ \\

    \theutterance \stepcounter{utterance}  
    & \multicolumn{4}{p{0.6\linewidth}}{
        \cellcolor[rgb]{0.561,0.737,0.561}{
            \makecell[{{p{\linewidth}}}]{
                \texttt{\tiny{[P1$\rangle$GM]}}
                \texttt{Good morning, everyone, and welcome to today's lecture on Japanese religious architecture and spirituality. As we explore the world of Shinto, it is crucial to understand that unlike many other religions, Shinto shrines are not merely buildings; they are sacred spaces designed to facilitate a direct connection between the human world and the divine.} \\
\\ 
\texttt{*The Core Concept of the Shrine*} \\
\texttt{At the heart of every Shinto shrine is the belief that kami, or spirits, reside within the structure itself. These spirits are not abstract concepts but are often associated with natural phenomena, ancestors, or historical figures. The architecture is intentionally simple and unadorned to avoid distracting from the presence of the kami. When you visit a site like the {[}address{]}, you are stepping into a space where the boundary between the mundane and the sacred is deliberately blurred, inviting visitors to purify themselves and offer respect.} \\
\\ 
\texttt{*Key Architectural Elements*} \\
\texttt{Inside the shrine, you will typically encounter specific structures that serve distinct functions. The most prominent is the *honden*, which houses the primary kami. It is often a small, enclosed room with a thatched roof, symbolizing the natural world. Surrounding this is the *haiden*, or main hall, where rituals are performed and where the kami is visually represented by a statue or a symbolic object. These spaces are arranged to guide the visitor through a process of purification and reverence, ensuring that one approaches the divine with a clean heart and mind.} \\
\\ 
\texttt{*The Atmosphere and Rituals*} \\
\texttt{Finally, the experience inside a Shinto shrine is defined by its atmosphere of quiet contemplation and active participation. Visitors often perform *handoyome*, a ritual of cleansing hands and mouth at the *chozuya* fountain before entering. Inside, the air is often filled with the scent of incense, and the sounds of chanting or the rustling of paper lanterns create a meditative environment. It is a place where history, nature, and spirituality converge, offering a profound sense of peace and connection to the past.} \\
\\ 
\texttt{Thank you for your attention, and I look forward to your questions after this session.} \\
            }
        }
    }
    & & \\ \\

    \theutterance \stepcounter{utterance}  
    & & & \multicolumn{2}{p{0.3\linewidth}}{
        \cellcolor[rgb]{0.9,0.9,0.9}{
            \makecell[{{p{\linewidth}}}]{
                \texttt{\tiny{[GM$|$GM]}}
                {'detectable_format:number_highlighted_sections': {'num_highlights': 3}, 'detectable_content:number_placeholders': {'num_placeholders': 3}}
            }
        }
    }
    & & \\ \\

    \theutterance \stepcounter{utterance}  
    & & & \multicolumn{2}{p{0.3\linewidth}}{
        \cellcolor[rgb]{0.9,0.9,0.9}{
            \makecell[{{p{\linewidth}}}]{
                \texttt{\tiny{[GM$|$GM]}}
                \texttt{game\_result = LOSE} \\
            }
        }
    }
    & & \\ \\

\end{supertabular}
}

\end{document}
